% !PRINTED 412
\weberpage{412}
%<<CONT
in $A$ there are also eight cyclic permutations of three figures each,
$\gamma = (0, 1, 2) \ldots$, and $A$ can be represented thus:
\[
A = B + B\gamma + B\gamma^2.
\]
Any one permutation $\alpha$ together with a permutation $\gamma$ may be
regarded as generators of the group. For if, say,
\[
\gamma = (0, 1, 2),\ \alpha_1 = (0, 1)(2, 3),
\]
then
\weq{1}{\alpha_1\gamma\alpha_1 = \gamma' = (0, 3, 1),}
and from $(0, 1, 2), (0, 3, 1)$ the whole group $A$ can be derived
(Bd. I, \S.\,153, 7.).

Likewise one may take $\gamma$ and $\gamma'$ as generating permutations
of $A$. In particular,
\weq{2}{\gamma'\gamma\gamma' = \alpha_1.}
We choose $\vartheta$ so that to the permutation $\gamma$ there corresponds a multiplication
$\varepsilon\vartheta$, in which, since $\gamma$ is of degree 3, $\varepsilon$ denotes a third root of unity.
 We set
\weq{3}{\vartheta = \frac{\varphi(x_0, x_1, x_2, x_3)}{\psi(x_0, x_1, x_2, x_3)},}
and understand by $\varphi, \psi$ two integral functions of the four
variables $x$ without common divisor.

From the identical equation
\[
\frac{\varphi_\gamma}{\psi_\gamma} = \varepsilon\,\frac{\varphi}{\psi}
\]
it then follows that
\weq{4}{\varphi_\gamma = \varepsilon_1\varphi,\ \psi_\gamma = \varepsilon_2\varphi} % [sic]
must hold, in which $\varepsilon_1, \varepsilon_2$ are likewise third roots of unity.

We apply one of the permutations $\alpha$ to $\vartheta$, and obtain
 from $\vartheta_\alpha = \chi(\vartheta)$:
\weq{5}{\frac{\varphi_\alpha}{\psi_\alpha} = \frac{a\varphi + b\psi}{c\varphi + d\psi}.}
Now since $\varphi, \psi$, hence also $a\varphi + b\psi$ and $c\varphi + d\psi$, and
likewise $\varphi_\alpha, \psi_\alpha$, are without common divisor, in the
identity (5) the numerator must equal the numerator and the denominator the
denominator, at least up to a constant factor.
This constant factor we can absorb into the constants $a, b, c, d$
and obtain
\[
\text{6)}\quad \varphi_\alpha = a\varphi + b\psi,\ \psi_\alpha = c\varphi + d\psi.
\]
